\section{Morphology breadth: aerial and humanoid}
\label{app:morphology}

The \S\ref{sec:onboarding} scoreboard covers serial arms and quadrupeds.
To test whether the \emph{same} pipeline (unchanged except for a
class-specific generation prompt and physics validator) generalizes to
qualitatively different dynamics, we onboard two further morphologies from
MuJoCo Menagerie: a free-flying quadrotor (Skydio~X2) and a bipedal humanoid
(Unitree~H1). Both are floating-base and open-loop unstable, so the
synthesized driver must contain a stabilizing feedback controller, not just
kinematics. We keep these two outside AutoAdapter-Bench proper (the eight
robots of \S\ref{sec:protocol}) and report them as breadth tests instead:
on each we run driver synthesis and the same-API CaP comparison, but not
the learned baselines (PPO, Diffusion Policy, OpenVLA) and seven-model
leaderboard that anchor the benchmark, so folding them into the benchmark count
would overstate how much head-to-head evidence backs them.

\noindent\textbf{Aerial (Skydio~X2).} The agent synthesizes a driver exposing
\texttt{takeoff}, \texttt{move\_to}, \texttt{hover}, \texttt{land}, and
\texttt{get\_base\_pose}; internally it implements a cascaded
position\,$\rightarrow$\,attitude\,$\rightarrow$\,rotor-thrust controller for
the underactuated 4-rotor base. It passes the aerial validation suite
($7/7$: build, \texttt{home}, \texttt{takeoff} to commanded altitude while upright, and
\texttt{move\_to} to a world target without flipping) on a Sonnet~4.5 agent.
On an 8-task aerial suite (4 simple + 4 hard; takeoff-to-altitude, hover,
point-to-point, climb, round-trip, diagonal, square patrol, tight hover;
physics-graded on post-step base pose with an upright constraint, and
hardened so that a do-nothing agent scores $0/8$: round-trips require real
movement, hovers require the commanded altitude, the patrol's waypoint
tolerance is below the square's center-to-corner distance), Auto-Adapter and
same-API CaP \emph{both} reach $100\%$ (Table~\ref{tab:aerial}, $N{=}5$;
filmstrips in Figure~\ref{fig:filmstrips2}).
Unlike contact-rich manipulation (\S\ref{sec:cap-portability}),
where the multi-turn loop confers an advantage,
point-to-point flight through the synthesized high-level API
(\texttt{takeoff}, \texttt{move\_to}) is solvable by a single-shot program, so
A-A and CaP tie. The aerial contribution is therefore the \emph{validated
flight controller itself} (a cascaded thrust stabilizer synthesized zero-shot
from the MJCF), not an agent-loop win.

\begin{table}[H]
\centering
\small
\begin{tabular}{lcc}
\toprule
\textbf{Aerial task} & \textbf{A-A} & \textbf{CaP} \\
\midrule
takeoff to 0.5\,m        & \checkmark & \checkmark \\
hover at 0.5\,m          & \checkmark & \checkmark \\
go to (0.4,\,0,\,0.5)    & \checkmark & \checkmark \\
climb to 1.0\,m          & \checkmark & \checkmark \\
out-and-return           & \checkmark & \checkmark \\
diagonal go-to           & \checkmark & \checkmark \\
square patrol (4 wpts)   & \checkmark & \checkmark \\
tight hover              & \checkmark & \checkmark \\
\midrule
\textbf{Physics \% ($N{=}5$)} & \textbf{100} & \textbf{100} \\
\bottomrule
\end{tabular}
\caption{Skydio~X2 aerial suite on the agent-synthesized drone driver
(Sonnet~4.5). \checkmark\ = $5/5$ trials pass. Both agent styles solve all
eight on the hardened suite (a do-nothing controller scores $0/8$):
point-to-point flight through the synthesized high-level API is single-shot
solvable, so Auto-Adapter and CaP tie; the contribution is the validated flight
driver, not the agent loop.}
\label{tab:aerial}
\end{table}

\noindent\textbf{Humanoid (Unitree~H1).} The 19-DoF H1 is a tall inverted
pendulum: open-loop it falls. The agent synthesizes a joint-space PD balance
controller that passes onboarding validation $7/7$: it holds a \emph{stable
stand} (torso $0.98$\,m, upright $1.00$ on the body-$z$ axis for the full
hold) and \emph{squat-recovers} without falling over (torso dips then returns to
$0.97$\,m at upright $0.99$). The result is reproducible under the benchmark
protocol (home-initialized, as in validation and replay): stand $4/4$ and
squat-recover $3/3$ across re-runs. Synthesizing a validated stand-and-squat
balance controller for a 19-DoF biped zero-shot from the MJCF is, to our
knowledge, the hardest morphology the pipeline handles. We report the headline H1 driver at the
onboarding level (the orchestrator steps physics directly for stand
and squat), since a pure \emph{stand} replay can be passed by doing
nothing (a homed biped stays up if physics is not advanced); the
cross-model sweep below instead grades H1 on \emph{active}
stand-and-squat tasks, where a squat must lower the body and return,
so a do-nothing controller fails. We also \emph{attempted} dynamic walking:
the generator was prompted for a closed-loop micro-step gait, validator-certified
for forward progress in the starting direction, staying upright, and
bounded sideways/turning drift (so falling forward, standing still, and spinning cannot pass).
It did not yield a stable validator-certified walk across attempts: adding the
walking objective even broke the otherwise-reliable stand. Dynamic
bipedal locomotion is thus the pipeline's hardest open case:
static balance and
squat are synthesizable zero-shot, dynamic gait is not (yet).

\noindent\textbf{Cross-model self-synthesis across morphologies.} The
\S\ref{sec:cap-portability} ``writes API'' result (only the three strongest
models synthesize a working SO-101 grasp driver) raises a question: is
self-synthesis difficulty a property of the \emph{model} or of the
\emph{morphology}? We ran the from-scratch pipeline for all seven LLMs on three
non-gripper morphologies (Go2 quadruped, Skydio~X2 quadrotor, Unitree~H1
humanoid), two attempts each, scoring a cell as a pass only if the synthesized
driver clears the class-specific validator and is not hardcoded (no skeleton
import, no literal leg-name constants). Table~\ref{tab:xmorph-synth} reports
reliability; Table~\ref{tab:xmorph-task} the task success of the passing
drivers. Two patterns hold. \emph{(i)~Reliability is morphology-ordered}: more
models clear aerial flight than locomotion than the humanoid biped, and only
the two strongest models in this sweep (Opus, Sonnet) synthesize
across more than one morphology while the open/small models (Nova, Ministral, Qwen) clear none.
\emph{(ii)~Task success of the drivers that do validate is also
morphology-ordered}: flight $0.82$--$1.00$ $\gtrsim$ locomotion $0.70$--$0.82$ $>$
humanoid $0.50$ $>$ SO-101 grasp $0.28$--$0.46$ (\S\ref{sec:cap-portability}).
Self-synthesis is thus jointly gated by model capability \emph{and}
morphology: in task success, writing a correct low-level controller
gets harder from smooth flight, through
stepping locomotion and slow biped balance, to
contact-heavy grasping; in synthesis reliability the biped is the
rarest to validate (Table~\ref{tab:xmorph-synth}).

\begin{table}[H]
\centering\small
\setlength{\tabcolsep}{3.5pt}
\begin{tabular}{lccccccc}
\toprule
 & \textbf{Op} & \textbf{So} & \textbf{Ha} & \textbf{No} & \textbf{DS} & \textbf{Mi} & \textbf{Qw} \\
\midrule
Skydio~X2 (fly)  & 2/2 & 2/2 & 1/2 & 0/2 & 1/2 & 0/2 & 0/2 \\
Go2 (walk)       & 1/2 & 1/2 & 0/2 & 0/2 & 0/2 & 0/2 & 0/2 \\
H1 (biped)       & 1/2 & 0/2 & 0/2 & 0/2 & 0/2 & 0/2 & 0/2 \\
\bottomrule
\end{tabular}
\caption{Cross-model self-synthesis \textbf{reliability} (validated,
non-hardcoded driver / 2 attempts). Op=Opus~4.8, So=Sonnet~4.6, Ha=Haiku~4.5,
No=Nova~Pro, DS=DeepSeek~V3.2, Mi=Ministral~8B, Qw=Qwen3~32B.}
\label{tab:xmorph-synth}
\end{table}

\begin{table}[H]
\centering\small
\setlength{\tabcolsep}{4pt}
\begin{tabular}{lcccc}
\toprule
 & \textbf{Opus} & \textbf{Sonnet} & \textbf{Haiku} & \textbf{DeepSeek} \\
\midrule
Skydio~X2 (fly)  & 1.00 & 0.95 & 1.00 & 0.82 \\
Go2 (walk)       & 0.70 & 0.82 & --   & --   \\
H1 (biped)       & 0.50 & --   & --   & --   \\
\bottomrule
\end{tabular}
\caption{\textbf{Task success} (physics, $N{=}5$, simple+hard) of the
self-synthesized drivers that passed validation; -- = no validated driver
(synthesis failed). Models that validated on no morphology (Nova, Ministral,
Qwen) are omitted.}
\label{tab:xmorph-task}
\end{table}
